\documentclass[11pt]{article}
\usepackage[margin=1.1in]{geometry}
\usepackage{amsmath,amssymb,amsthm}
\usepackage{booktabs}
\usepackage{enumitem}
\newtheorem{proposition}{Proposition}
\newtheorem{assumption}{Assumption}
\newtheorem{definition}{Definition}
\newtheorem{remark}{Remark}
\newcommand{\E}{\mathbb{E}}
\newcommand{\R}{\mathbb{R}}
\newcommand{\enc}{\phi_\theta}
\newcommand{\head}{g_\theta}
\newcommand{\pol}{f_\theta}

\title{Why Mean Regression Fails Behavior Cloning:\\
A Unified Gradient- and Representation-Level Account\\
of the Regression--MIP Gap}
\author{Theory note draft}
\date{July 2026}

\begin{document}
\maketitle

\begin{abstract}
We give a unified account of why $L_2$ behavior-cloning regression underperforms
two-view denoising objectives (MIP) on both human-demonstrated and scripted
manipulation data, from the joint viewpoint of \emph{gradient dynamics} and
\emph{representation dynamics}. The organizing object is the \emph{realizability}
of the training target given the network input. For state-only regression the
target contains an unrealizable component --- cross-demonstration disagreement on
human data, near-degenerate label micro-differences on scripted data --- and the
two observed failure modes are the two ways this component can damage training:
(I) when it is heavy-tailed at macroscopic scale, it captures the batch gradient
of the shared encoder and destroys previously learned structure (optimization
pathology); (II) when it is microscopic and degenerate, it withdraws the
incentive to keep distinct states distinct, and the encoder folds (representation
pathology). The anchor view of MIP makes every training target realizable, which
simultaneously removes the noise gradient of (I) and creates a distinction
incentive of relative strength $(\delta/\sigma)^2$ against (II). On the
regression side, (I) is replaceable by a bounded-influence likelihood with
learned state-conditional scale (Student-$t$), and the conditioning deficit is
replaceable by observation history; (II) has no known loss-side replacement.
Every proposition is paired with a measured quantity from our experiments. The
note does not claim distributional multimodality is the operative difficulty;
the measured disagreement is timing jitter around a single strategy, and the
analysis proceeds entirely at the level of training signals.
\end{abstract}

\section{Setup}

Let $s \in \mathcal{S} \subset \R^{d_s}$ be an observation (possibly a window of
past frames) and $a \in \R^{d_a}$ an action chunk. Demonstrations supply pairs
$(s_i, a_i)$ with
\begin{equation}
a_i = \mu(s_i) + \eta_i, \qquad \E[\eta \mid s] = 0 ,
\end{equation}
where $\mu$ is the state-conditional mean and $\eta$ collects the component of
the action that is \emph{not a function of the input}: across demonstrations,
operators at indistinguishable states act at different times, speeds, and grasp
instants. The policy is $\pol = \head \circ \enc$ with a shared encoder $\enc$
and action head $\head$; both views of the analysis concern the gradient
received by the shared parameters of $\enc$.

\begin{definition}[Realizable component]
Given a network input $x$ (which may be $s$, an observation window, or an
augmented input), the target $a$ decomposes as $a = m^\star(x) + \eta(x)$ with
$m^\star(x) = \E[a \mid x]$. We call $\eta(x)$ the \emph{unrealizable component}
relative to $x$: no measurable function of $x$ can reduce it.
\end{definition}

Two empirically measured data regimes anchor the analysis.

\begin{assumption}[Human regime]\label{as:human}
$\eta \mid s$ is heavy-tailed with tail index $\nu \approx 2$ and strongly
state-dependent scale: $\mathrm{Var}(\eta \mid s)$ varies by orders of magnitude
across task phases, and the top-decile states by conditional scale carry a
disagreement dominated by discrete timing events (gripper open/close instants,
pause-versus-move). Measured: tail energy of the top decile is $85\%$ gripper
channel; kurtosis of the per-sample deviation $\approx 10.6$ (tool-hang),
$8.1$ (square-mh).
\end{assumption}

\begin{assumption}[Scripted regime]\label{as:script}
$\eta \approx 0$, but the labels are \emph{degenerate}: there exist large state
regions $R$ on which labels differ only microscopically, $a(s) = \bar a + \delta
u(s)$ with $\|\delta u\| \ll$ the typical action scale, while closed-loop
correctness requires distinguishing states inside $R$.
\end{assumption}

\section{Gradient dynamics under $L_2$ (human regime)}

The per-sample loss $\ell_i = \|\pol(s_i) - a_i\|^2$ contributes to the encoder
the gradient
\begin{equation}
g_i \;=\; -2\, J_i^\top r_i, \qquad
J_i = \frac{\partial \pol}{\partial \enc}(s_i), \quad r_i = a_i - \pol(s_i).
\label{eq:persample}
\end{equation}
Once the realizable part is fit, $r_i \to \eta_i$: \emph{the residual, and hence
the per-sample encoder gradient, is the unrealizable component itself}.

\begin{proposition}[Gradient capture and churn]\label{prop:capture}
Under Assumption~\ref{as:human} with tail index $\nu \le 2$,
$\E\|\eta\|^2 = \infty$ (Student-$t$ boundary case), and for a batch
$B$ the noise part of the batch gradient $\sum_{i \in B} J_i^\top \eta_i$ is
dominated by its top order statistics: the ratio of the largest summand to the
sum does not vanish as $|B| \to \infty$ (regular variation), and the
signal-to-noise ratio of the batch mean gradient tends to zero. Consequently the
shared encoder performs a heavy-tailed random walk driven by irreducible
residuals, which degrades components of $\enc$ that carry already-learned
structure.
\end{proposition}

\emph{Derivation sketch.} For regularly varying $\|\eta\|$ with index
$\nu \le 2$, sums and maxima are of the same order
(single-big-jump principle); the empirical mean does not concentrate. Since
$\eta \perp$ the learnable signal, the noise gradient is an unbiased zero-mean
perturbation of unbounded variance applied to \emph{shared} parameters, i.e.\ a
random reparameterization of features used elsewhere. \hfill$\square$

\noindent\textbf{Measured counterparts.} Top-$10\%$ of samples carry
$74\%$--$77\%$ of the encoder gradient under $L_2$; the influence distribution's
kurtosis grows $36 \to 4041$ over training (spike churn); behavioral
acquire-then-unlearn (partial insertion servo at $37$k steps, lost by $300$k,
SR $43 \to 31$); fit on the SR-critical low-disagreement samples is up to
$9\times$ worse than under the robust loss at matched steps.

\begin{proposition}[Bounded influence restores the signal]\label{prop:robust}
Replace $\ell_i$ by the Student-$t$ negative log-likelihood with
state-conditional scale $\sigma(s)$,
$\rho(r) = \tfrac{\nu+1}{2}\log\!\big(1 + \|r\|^2/(\nu\sigma^2 D)\big)$.
Its influence function
$\psi(r) = (\nu+1)\, r / (\nu\sigma^2 + \|r\|^2/D)$
is bounded, so the noise-gradient has finite variance and the batch SNR is
restored; and with $\sigma(s)$ fit by maximum likelihood the per-sample weight
is asymptotically inverse to the conditional scale --- the efficient reweighting
under the assumed family. The predicted influence ratio between top-decile
disagreement samples and typical samples is
$\psi(\eta_{\mathrm{tail}})/\psi(\eta_{\mathrm{typ}}) < 1$ once
$\sigma(s)$ calibrates.
\end{proposition}

\noindent\textbf{Measured counterparts.} Own-view top-$1\%$ share $5$--$9\%$
late (versus $59\%$ for $L_2$), kurtosis $9.5$--$23$; tail/typical influence
$0.60$; the calibration transient is visible ($\approx 50$k steps of
$L_2$-like behavior before $\sigma$ fits); SR $31$--$43 \to 75$--$88$. Hard
removal of the same samples (trimming, kNN smoothing) \emph{hurts} --- the
information is bounded, not deleted --- consistent with the estimator view and
inconsistent with a ``noise filtering'' view.

\begin{proposition}[The anchor removes the unrealizable component]\label{prop:anchor}
The MIP anchor view supplies input $x = (s,\, a + \sigma\varepsilon,\, t^\star)$
with target $a$. Then $a$ is realizable given $x$ up to the Bayes error of
denoising, which the identity-plus-correction map attains; residuals vanish for
\emph{every} sample, including the disagreement tail. Hence the anchor view
contributes no irreducible noise gradient, and no reweighting is needed:
$g(\mathrm{tail})/g(\mathrm{typ}) \approx 1$.
\end{proposition}

\noindent\textbf{Measured counterparts.} MIP own-view tail ratios $0.85$--$1.04$
and top-$1\%$ shares $3$--$17\%$ throughout training; evaluating the
\emph{anchor-free} readout of an anchor-trained network shows residuals
concentrated precisely on the tail (ratios $3.9$--$4.1$) --- the network never
learned those targets from the state; the anchor carried them.

\begin{remark}[Conditioning as an alternative route to realizability]
Realizability can also be improved on the input side: if part of $\eta$ is
predictable from longer observation history $h$, then
$\mathrm{Var}(\eta \mid s, h) < \mathrm{Var}(\eta \mid s)$ and the unrealizable
component shrinks before the loss ever acts. Measured: extending the window from
$2$ to $8$ frames cuts the tail-sample fitting error $4$--$5\times$ for
\emph{both} losses early in training, while leaving the learnable-majority fit
nearly unchanged; the loss and history factors act on complementary terms, and
success rates combine superadditively ($12/28/35/60$ across the $2\times2$).
\end{remark}

\begin{remark}[Estimator location]
Where $\eta \mid s$ is a mixture over timing (moving versus paused), the
conditional mean is itself a biased action: on transport its speed is
$\approx 0.5\times$ the demonstrated speed in closed loop. No reweighting moves
a location estimate off the mean; a quantile relocation along the demonstrated
direction (pinball loss on the parallel speed component, $\tau > 1/2$) does,
and only there. This is a second, logically independent defect of the mean.
\end{remark}

\section{Representation dynamics under label degeneracy (scripted regime)}

\begin{proposition}[No distinction incentive]\label{prop:fold}
Under Assumption~\ref{as:script}, let $\enc$ be \emph{collapsed on $R$}
($\enc(s) = \bar\phi$ for $s \in R$). The excess $L_2$ risk of the collapsed
encoder is $O(\|\delta\|^2)$, and the loss gradient toward separating $R$'s
states scales as $O(\|\delta\|^2)$ as well --- for microscopic $\delta$ it is
dominated by any competing compression pressure (weight decay, simplicity bias).
The encoder therefore folds: states requiring different downstream behavior
share a representation cell. The train loss cost is $O(\|\delta\|^2)$, but the
closed-loop cost is $O(1)$: at deployment, neighborhoods are misassigned and the
policy reads out the wrong region's behavior.
\end{proposition}

\noindent\textbf{Measured counterparts.} Distinction sensitivity of $L_2$ and
Cauchy on scripted data $0.95$--$0.98$ (blind); settle-phase embedding
neighborhoods $4\%$ pure; deployed failure is an attenuated readout of
wrong-phase actions (cosine $+0.97$ to the wrong region's labels at half
amplitude). Note the loss-side lever of Section~2 is inert here: robust
reweighting acts on residual \emph{scale}, and all residuals are microscopic.

\begin{proposition}[The anchor creates a distinction incentive of relative
strength $(\delta/\sigma)^2$]\label{prop:incentive}
Consider two states $s_1, s_2 \in R$ with labels $\bar a \pm \delta/2$ and the
anchor view with noise level $\sigma \gg \|\delta\|$. If $\enc$ is collapsed,
the denoiser must estimate the label from the anchor alone; the two-point
mixture posterior yields excess denoising risk
\begin{equation}
\Delta \;=\; \frac{\|\delta\|^2}{4}\,
\Big(1 - \E\tanh^2\!\big(\tfrac{\delta^\top(\sigma\varepsilon)}{2\sigma^2} +
\tfrac{\|\delta\|^2}{4\sigma^2}\big)\Big)
\;\approx\; \frac{\|\delta\|^2}{4}\Big(1 - O\big(\tfrac{\|\delta\|^2}{\sigma^2}\big)\Big),
\end{equation}
which a separated encoder reduces to zero. Normalizing by the view's overall
loss scale $\sigma^2$, the \emph{relative} incentive to separate is
$\Theta\!\big((\|\delta\|/\sigma)^2\big)$ per unit of anchor-view weight ---
whereas for $L_2$ the relative incentive is $\Theta(\|\delta\|^2)$ against an
$O(1)$ loss scale. The anchor thus amplifies microscopic label distinctions
into a first-order representation pressure, with an explicit dose control
$\sigma$.
\end{proposition}

\noindent\textbf{Measured counterparts.} Distinction sensitivity of the anchor
view $13\times$ that of $L_2$ on scripted data, with measured dose--response
$(\delta/\sigma)^2$; smaller $\sigma$ increases the incentive (until anchors
under-mix). We know of no loss reweighting that reproduces this: any
$\rho(\|r\|)$ evaluated at $\|r\| = O(\delta)$ gives an incentive vanishing with
$\delta$.

\section{The unified statement}

\begin{center}
\begin{tabular}{@{}lll@{}}
\toprule
 & \textbf{Human regime} & \textbf{Scripted regime} \\
\midrule
unrealizable component & macroscopic, heavy-tailed & microscopic, degenerate \\
pathology & gradient capture, churn (Prop.~\ref{prop:capture}) & representation fold (Prop.~\ref{prop:fold}) \\
level & optimization & representation \\
MIP's fix & realizability (Prop.~\ref{prop:anchor}) & incentive $(\delta/\sigma)^2$ (Prop.~\ref{prop:incentive}) \\
regression-side fix & Student-$t$ + history (+ quantile) & \emph{open} \\
measured signature & top-share, kurtosis, tail ratio & sensitivity, neighborhood purity \\
\bottomrule
\end{tabular}
\end{center}

Both regimes instantiate one principle: \emph{$L_2$ regression charges the
network for the component of the target it cannot realize, and the damage is
determined by the scale distribution of that component.} Heavy macroscopic
scale damages optimization; degenerate microscopic scale damages
representation. The anchor view is a single mechanism that neutralizes both
(realizability removes the charge; the anchor-conditioned Bayes problem
re-prices microscopic distinctions at $(\delta/\sigma)^2$). On the regression
side the first regime is fully replaceable by matched training signals; the
second is not, and this asymmetry is exactly where a denoising-class objective
is irreplaceable.

\section{Relation to Jacobian-regularization accounts}

A complementary hypothesis holds that denoising objectives implicitly regularize
the conditioning Jacobian of the recovered clean map (schedule coefficient
$\kappa_t$). Two of our measurements bound its role in the present setting.
First, for the $x_0$-parameterization used by MIP, $\kappa_t = 1$ and no
shrinkage is implied; consistently, measured boundary effective-Lipschitz
constants of MIP and regression policies coincide ($2.6$--$3.7$) while their
success rates differ by $40$ points. Second, enforcing output flatness directly
(finite-difference penalties at the deficit radius, normal-direction jitter,
shell-derivative penalties) degrades or destroys the policy at every dose tried,
while the successful interventions leave the Jacobian unchanged. In this domain
the operative channels of the denoising signal are the two identified above ---
gradient realizability and representation incentive --- not clean-map
smoothness. The Jacobian account may govern other regimes (e.g.\ genuinely
low-variance conditionals with $\kappa_t < 1$ parameterizations); the two
accounts are logically independent and empirically separable by the probes used
here.

\section{Design implications}

\begin{enumerate}[leftmargin=1.4em]
\item \textbf{Diagnose the unrealizable component first}: its tail index,
alignment with a data-defined minority, channel decomposition, and
history-resolvability determine which fix applies.
\item \textbf{Human-style data}: bounded-influence likelihood with learned
state-conditional scale (Student-$t$, $\nu$ set to the measured tail index);
observation history to the resolvability plateau; quantile relocation only on
coordinates where the mean is demonstrably biased.
\item \textbf{Scripted-style data}: no loss reweighting suffices; a
distinction incentive is required (anchor view with $\sigma$ chosen by the
$(\delta/\sigma)^2$ trade-off, or data-side diversification).
\item \textbf{What not to do}: hard sample removal (deletes bounded
information), isotropic input noise with fixed labels (destroys closed-loop
gain), output-flatness penalties (attack magnitude, not direction), inverse-residual
channel weighting (re-concentrates on noise floors).
\end{enumerate}

\section{Limitations}

The propositions are stated at the level of estimator and incentive structure,
not end-to-end learning dynamics; the churn mechanism in
Proposition~\ref{prop:capture} is supported by measured stream statistics and
behavioral unlearning rather than a convergence theorem. The $(\delta/\sigma)^2$
computation treats a two-point mixture with a collapsed encoder; the constant in
front depends on architecture. Success-rate arithmetic (which points a fix
recovers) is task-dependent and is reported empirically in the companion
experimental tables.

\end{document}
